\section{Challenges}
\label{sec:challenges}

State legislative redistricting poses two computational challenges that do not
arise at the congressional level: the high-$k$ problem and the resolution problem.

\subsection{The High-$k$ Problem}
\label{sec:challenges:highk}

Congressional redistricting requires partitioning a state into at most 52 parts
(California, 2020). The METIS multilevel $k$-way partitioner
\citep{karypis1998fast} handles this comfortably at census-tract resolution:
California's $\sim\!9{,}000$ tracts divided into 52 districts gives
$\sim\!173$ tracts per district on average, well above the minimum needed
for meaningful boundary optimisation.

State legislative redistricting changes this picture dramatically. New Hampshire's
400-seat house must be drawn from $\approx 320$ census tracts---fewer tracts
than districts. Wyoming's 60-seat house covers $\approx 140$ tracts, giving a
ratio of 0.43. Even Washington state's 98-seat house has a ratio of
$98/1{,}458 \approx 0.067$, above the threshold where tract resolution degrades
compactness optimisation.

The high-$k$ problem has two manifestations:
\begin{enumerate}
  \item \textbf{Degeneracy}: when $k > n_{\mathrm{tracts}}$, METIS cannot produce
    $k$ non-empty parts---some districts must consist of fractional tracts.
    The pipeline switches to block-group or block resolution to restore
    $k \ll n_{\mathrm{units}}$.
  \item \textbf{Coarseness}: when $k/n_{\mathrm{tracts}}$ approaches $0.05$,
    each district contains on average only 20 tracts. The boundary between
    adjacent districts is defined by only a few shared edges, making
    Polsby-Popper scores artificially high and edge-cut optimisation less
    meaningful. Resolution improvement restores the geometric richness of the
    boundary set.
\end{enumerate}

\subsection{The Resolution Problem}
\label{sec:challenges:resolution}

The \texttt{redist} pipeline currently supports three resolution levels:
\begin{itemize}
  \item \textbf{Census tract}: $\sim\!73{,}000$ nationally, 1,200--8,000
    residents each. Standard for congressional redistricting.
  \item \textbf{Census block group}: $\sim\!220{,}000$ nationally, 600--3,000
    residents each. Needed for high-$k$ state chambers.
  \item \textbf{Census block}: $\sim\!8{,}100{,}000$ nationally, 0--600
    residents each. Required only for the most extreme cases (NH, WY).
\end{itemize}

Moving from tract to block-group resolution multiplies the number of geographic
units by $\sim\!3\times$ and increases the adjacency graph size by a similar
factor. This increases runtime by approximately 3--4$\times$ per state.
Building block-group adjacency data requires a separate preprocessing step:
\begin{verbatim}
python scripts/data/geography/build_unit_adjacency.py \
  --resolution block_group --year 2020
\end{verbatim}
This is a one-time cost per census year; the resulting \texttt{.adj.bin} files
are cached and reused across all runs at that resolution.

\subsection{The Resolution Selection Rule}
\label{sec:challenges:rule}

We derive the resolution selection rule empirically and theoretically in
companion paper F.3 \citep{dellaLibera2026stateLegResolution}. The rule is:

\begin{definition}[Resolution selection rule]
  Let $k$ be the number of districts and $n_{\mathrm{tracts}}$ be the number of
  census tracts in the state. Use block-group resolution if and only if
  \[
    \frac{k}{n_{\mathrm{tracts}}} > 0.05.
  \]
\end{definition}

Intuitively: each district should contain at least 20 base geographic units
on average ($k \cdot 20 \leq n_{\mathrm{tracts}}$, i.e., $k/n \leq 0.05$).
Below this threshold, the boundary set is rich enough for meaningful compactness
optimisation; above it, resolution upgrade is necessary.

In practice, most state house chambers have $k/n < 0.05$ and can use tract
resolution. The exceptions are principally small states with large lower chambers:
New Hampshire, Wyoming, Vermont, and South Dakota. Washington state sits at the
boundary ($k/n \approx 0.067$) and benefits meaningfully from block-group resolution.
